\documentclass[11pt]{article}
\usepackage[margin=1in]{geometry}
\usepackage{amsmath,amssymb,amsthm}
\usepackage[round]{natbib}
\usepackage{hyperref}
\usepackage{microtype}

\newtheorem{theorem}{Theorem}
\newtheorem{proposition}{Proposition}
\newtheorem{lemma}{Lemma}
\newtheorem{corollary}{Corollary}
\newtheorem{remark}{Remark}

\title{Betting Against a Conformal Predictor:\\
a Parimutuel Account of the Information Gap}
\author{Peter Cotton}
\date{Draft --- companion note}

\begin{document}
\maketitle

\begin{abstract}
The companion paper shows that the log-score regret of a single-shape conformal predictor to the
conditional oracle is the mutual information $I(R;X)$ between the residual and the input. Here we
rederive that quantity from a betting mechanism rather than from the score. Treat the predictor as
the crowd in a parimutuel pool on the residual: bettors put money in, and the pot is split among
the winners in proportion to their stake on the realised outcome. In the continuous limit the
payoff is the ratio of the bettor's density to the crowd's. An entrant who knows only the marginal
breaks even, which is marginal coverage stated as wealth. An entrant who conditions on $X$ grows
his bankroll at rate exactly $I(R;X)$. The gap is the rent. This is not a metaphor: the
nearest-the-pin pool of the microprediction platform, and the continuous density version run in the
MidOne contest, are this mechanism, and a conformal predictor is the entrant that prices the pool
flat in $X$. We then give two ways to measure the rent on a fitted predictor: a static lower bound
from distance covariance, and a sequential e-process whose growth rate estimates $I(R;X)$ and which
is an anytime-valid test for conditional miscoverage.
\end{abstract}

\section{Setup}

A point predictor leaves a residual $R$, and conformalization re-levels its marginal law. Write
$U=G(R)$ for the probability integral transform, $G$ the marginal CDF of $R$; $U$ is marginally
uniform on $[0,1]$ whatever the model. Let $g(u\mid x)$ be the conditional density of $U$ given
$X=x$. The companion paper \citep{cotton_marginally} identifies the wasted log-score of a
single-shape conformal system as
\begin{equation}\label{eq:gap}
\Delta \;=\; \mathbb E_X\,\mathrm{KL}\!\big(P_{R\mid X}\,\big\|\,G\big)
\;=\; I(R;X) \;=\; I(U;X) \;=\; \mathbb E_X\!\int_0^1 g(u\mid X)\log g(u\mid X)\,\mathrm du .
\end{equation}
The middle equalities are the invariance of mutual information under the deterministic, a.s.
invertible map $R\mapsto G(R)$, and the fact that the marginal of $U$ is uniform, so the
information is the negative conditional differential entropy of $U$. We take \eqref{eq:gap} as
given and ask what it is, operationally. The answer is a betting rate.

\section{The continuous parimutuel}

A parimutuel sets no odds. Everyone's money goes into one pool; once the outcome is known the whole
pool is paid to the holders of winning tickets, divided in proportion to their winning stake. With a
continuous outcome the winning ``ticket'' is a density, and the payoff has a clean form.

\begin{lemma}[Pool payoff]\label{lem:payoff}
Normalise the pool to $1$ and let the crowd's aggregate stake have density $q$ on the outcome space,
$\int q=1$. An infinitesimal, price-taking entrant who stakes unit wealth with density $b$,
$\int b=1$, and takes no part in moving $q$, holds after outcome $u$ the wealth
\begin{equation}
W(u)\;=\;\frac{b(u)}{q(u)} .
\end{equation}
\end{lemma}

\begin{proof}
Bin the outcome axis at width $\delta$. The crowd stakes $q(u)\delta$ of the unit pool on the bin at
$u$; the entrant stakes $b(u)\delta$. If the outcome lands in that bin, the entire pool is split
among its stakeholders in proportion to stake, so the per-unit payoff there is $1/(q(u)\delta)$ (the
entrant is infinitesimal, so the crowd's stake is the whole bin), and the entrant collects
$b(u)\delta \cdot 1/(q(u)\delta)=b(u)/q(u)$. The bin width cancels, so the limit is well posed and
equals the Radon--Nikodym derivative $\mathrm db/\mathrm dq$.
\end{proof}

The cancellation is the point. A continuous lottery on a point outcome looks ill defined --- every
exact value has probability zero --- but pool-splitting divides two vanishing quantities, and the
ratio survives. No reference measure and no assumed ``fair odds'' enter.

\begin{lemma}[Log-optimal growth]\label{lem:growth}
Let the outcome be drawn from the true law $p$. The expected log-growth of an entrant staking $b$
against a crowd staking $q$ is
\begin{equation}
G(b)=\mathbb E_{U\sim p}\!\left[\log\frac{b(U)}{q(U)}\right]
=\mathrm{KL}(p\|q)-\mathrm{KL}(p\|b)\;\le\;\mathrm{KL}(p\|q),
\end{equation}
with equality iff $b=p$ a.e. The log-optimal stake is the truth, and the optimal growth rate is
$\mathrm{KL}(p\|q)$.
\end{lemma}

\begin{proof}
Add and subtract $\log p$: $\int p\log(b/q)=\int p\log(p/q)-\int p\log(p/b)$. The first term is
$\mathrm{KL}(p\|q)$, the second is $\mathrm{KL}(p\|b)\ge 0$ by Gibbs' inequality, zero iff $b=p$.
\end{proof}

Lemmas \ref{lem:payoff} and \ref{lem:growth} are the parimutuel statement of the Kelly--Breiman
log-optimality principle \citep{kelly1956,cover2006}; the growth rate $\mathrm{KL}(p\|q)$ is the
information the crowd's price $q$ is missing about the truth $p$.

\section{The conformal lottery}

Put the pool on the rank scale. Conformalization makes $U$ marginally uniform, so the conformal
predictor is exactly the crowd that prices the rank pool flat and uses nothing about $X$:
\begin{equation}
q(u)\equiv 1 \quad\text{on } [0,1].
\end{equation}
Two entrants step up.

\begin{theorem}[Parimutuel rent of conformal]\label{thm:rent}
Against the conformal crowd $q\equiv 1$:
\begin{enumerate}
\item[(i)] an entrant who knows only the marginal stakes $b\equiv 1$ and grows at rate $0$;
\item[(ii)] an entrant who observes $X$ and stakes the conditional $b=g(\cdot\mid X)$ grows, per
round, at rate
\begin{equation}
G^\star=\mathbb E_X\,\mathrm{KL}\!\big(g(\cdot\mid X)\,\big\|\,1\big)=I(R;X).
\end{equation}
\end{enumerate}
With a track take $\tau\in[0,1)$ (payoff $(1-\tau)/q$) the informed entrant grows at
$I(R;X)+\log(1-\tau)$, positive iff $I(R;X)>-\log(1-\tau)\approx\tau$.
\end{theorem}

\begin{proof}
Apply Lemma \ref{lem:payoff} with $q\equiv 1$, then Lemma \ref{lem:growth}. For (i) the entrant's
true law is the marginal, uniform, so $b=q=1$ and $G=\mathrm{KL}(1\|1)=0$. For (ii) condition on
$X=x$: the entrant's correct law is $g(\cdot\mid x)$, so his conditional growth is
$\mathrm{KL}(g(\cdot\mid x)\|1)$; averaging over the $X$ he sees each round gives
$\mathbb E_X\,\mathrm{KL}(g(\cdot\mid X)\|1)=I(U;X)=I(R;X)$ by \eqref{eq:gap}. The take multiplies
every payoff by $(1-\tau)$, adding $\log(1-\tau)$ to the rate.
\end{proof}

So the certificate and the leak are two readings of the same pool. Marginal coverage says the
lottery is fair to anyone who knows only the average; part (i) is that fairness, as a break-even.
Part (ii) is what the certificate cannot see: a side-informed entrant compounds at $I(R;X)$,
bankroll growing like $e^{tI(R;X)}$, and no amount of re-leveling closes the gap, because re-leveling
acts on the marginal where the leak is invisible. Conformal re-leveling takes no rake ($\tau=0$), so
any conditional information at all is pure profit. This is the lottery paradox in one line: a
provably fair lottery that a player with side information beats with certainty.

\section{The mechanism is real}

The continuous parimutuel is not a thought experiment. It is the reward rule of a deployed
prediction market.

The microprediction platform \citep{cotton_micro2022} asks each entrant for $225$ Monte Carlo
samples of the next value of a live scalar stream. When the value $z$ arrives, a stateless
parimutuel splits the pot toward the entries nearest $z$ --- nearest-the-pin. Smoothed, an entry's
reward is its sample density at $z$ relative to the field's aggregate sample density at $z$, which
is the payoff $b(z)/q(z)$ of Lemma \ref{lem:payoff} with the densities estimated from samples. The
incentive is log-optimal by Lemma \ref{lem:growth}: an entrant maximises expected log-reward by
sampling from his true conditional, $b=p$, so the aggregate crowd price $q$ is driven toward the
truth.

The continuous version drops the samples and prices a density directly. In the MidOne contest run by
Crunch Labs (CrunchDAO), entrants submit an explicit predictive density --- a mixture, in the
conventions of the \texttt{density} package \citep{micro_density}, evaluated pointwise --- and the
wealth change splits the pot in proportion to the density the entrant ascribes to the realised
outcome against the density the field ascribes to it, again $b(z)/q(z)$. MidOne scores attacks on
small departures from martingality, so the priced object is a residual density: the $R$ of this
note.

A conformal predictor, dropped into either contest, is the entrant who prices the residual pool flat
in $X$ --- correct on the margin, indifferent to the case in front of it. Theorem \ref{thm:rent}
then says any entrant who conditions on $X$ holds a guaranteed log-growth rate $I(R;X)$ against it.
In these markets the information gap is not a figure of speech. It is the realised bankroll of the
better-informed competitor.

\section{Measuring the rent}

The rent is defined through a log-score, so it is not read off data directly. Two routes estimate
it on a fitted predictor: a cheap lower bound, and an exact sequential estimate.

\subsection{A static lower bound from distance covariance}

The companion diagnostic \citep{cotton_signed} tests $U\perp X$ with distance covariance, which is
zero iff the gap is zero. The same statistic bounds the gap from below.

\begin{lemma}[Certified lower bound]\label{lem:lower}
Let a bounded characteristic kernel with $k\le K$ on the product space induce
$\operatorname{dCov}(U,X)=\operatorname{HSIC}(U,X)$ \citep{sejdinovic2013}. Then
\begin{equation}
I(R;X)\;\ge\;\frac{\operatorname{HSIC}(U,X)}{2K}.
\end{equation}
\end{lemma}

\begin{proof}
HSIC$^{1/2}$ is the maximum mean discrepancy $\operatorname{MMD}(P_{U,X},P_U\otimes P_X)$, an
integral probability metric over the unit RKHS ball. Each such witness $f$ has
$\|f\|_\infty\le\sqrt K\|f\|_{\mathcal H}\le\sqrt K$, so
$\operatorname{MMD}\le 2\sqrt K\,\mathrm{TV}(P_{U,X},P_U\otimes P_X)$. Pinsker's inequality
\citep{cover2006} gives $\mathrm{KL}\ge 2\,\mathrm{TV}^2$, and the KL of the joint against the
product of marginals is $I(U;X)=I(R;X)$. Combining, $I(R;X)\ge 2(\operatorname{MMD}/2\sqrt K)^2=
\operatorname{HSIC}/(2K)$.
\end{proof}

A measured distance covariance therefore certifies a minimum wasted log-score, and by Theorem
\ref{thm:rent} a minimum extractable rent. The matching upper bound does not exist with a fixed
kernel: MMD metrizes weak convergence \citep{simongabriel2023} whereas KL does not, so a sequence of
laws can have HSIC$\to 0$ with $I(R;X)\to\infty$. Under a bounded conditional density ratio
$g(u\mid x)\le M$ one has the crude $I(R;X)\le\log M$, refinable through $\chi^2$
\citep{wangtay2022}, but the exact value needs a different object.

\subsection{A sequential estimate: the e-process}

Run the informed entrant and watch the bankroll. At round $t$, having seen
$(U_1,X_1),\dots,(U_{t-1},X_{t-1})$, choose a betting density $b_t(\cdot\mid x)\ge 0$ with
$\int_0^1 b_t(u\mid x)\,\mathrm du=1$ for every $x$, and set
\begin{equation}
W_t=\prod_{s\le t} b_s(U_s\mid X_s),\qquad W_0=1.
\end{equation}

\begin{proposition}[Anytime-valid test and consistent rent estimate]\label{prop:eprocess}
\quad
\begin{enumerate}
\item[(i)] Under $H_0:U\perp X$ (the conformal price is correct, $I(R;X)=0$), $(W_t)$ is a
non-negative martingale with $\mathbb E\,W_t=1$, so by Ville's inequality
$\mathbb P(\sup_t W_t\ge 1/\alpha)\le\alpha$. Rejecting when $W_t\ge 1/\alpha$ is an anytime-valid,
level-$\alpha$ test of conditional miscoverage.
\item[(ii)] If the entrant learns the conditional, $b_t\to g(\cdot\mid x)$, then
$\tfrac1t\log W_t\to I(R;X)$ almost surely.
\end{enumerate}
\end{proposition}

\begin{proof}
(i) Under $H_0$, $U_t\mid X_t$ is uniform and independent of the past, so
$\mathbb E[b_t(U_t\mid X_t)\mid\mathcal F_{t-1}]=\int_0^1 b_t(u\mid X_t)\,\mathrm du=1$; the product
of conditionally mean-one factors is a martingale with mean $1$, and Ville's inequality applies.
(ii) $\tfrac1t\log W_t=\tfrac1t\sum_s\log b_s(U_s\mid X_s)\to\mathbb E[\log g(U\mid X)]$ by the law
of large numbers once $b_s\to g$, and $\mathbb E[\log g(U\mid X)]=\mathbb E_X\int g(u\mid X)\log
g(u\mid X)\,\mathrm du=I(R;X)$ by \eqref{eq:gap}.
\end{proof}

This is the betting story made operational. The wealth process is at once the test --- it crosses
$1/\alpha$ only if coverage is conditionally uneven --- and the estimator, its growth rate
converging to the rent the static bound could only lower. It is the log-optimal e-variable of safe
testing \citep{gruenwald2024,vovkwang2021} for the null $U\perp X$, and the sequential, anytime
sibling of the batch distance-covariance test. Choosing $b_t$ is online conditional density
estimation; any consistent learner gives the convergence, and a poor one still keeps the type-I
guarantee in (i), since validity needs only $\int b_t=1$.

A measurement caveat, shared with the companion diagnostic. The transform $U=G(R)$ uses the true
marginal $G$; with the empirical rank the null is exact only in the leave-one-out or full-conformal
sense, which is what we use to keep $U_t\mid X_t$ uniform under $H_0$.

\section{Scope}

This detects; it does not certify. By the distribution-free no-go results
\citep{leiwasserman2014,foygelbarber2021,vovk2012} no procedure can prove conditional coverage from
data alone, so a quiet e-process or a small distance covariance means ``no rent found at this
power,'' never conditional validity. What the betting view adds is a positive reading of the other
side: when the gap is there, it is money, and it is extractable by anyone who conditions on $X$.
Closing it is a modelling job --- a sharper conditional spread, not a finer threshold --- and the
rent is the exact bounty for doing it.

\section{Related work}

The identity $\Delta=I(R;X)$ is the companion paper \citep{cotton_marginally}; the
distance-covariance diagnostic and the signed de Finetti reading are \citep{cotton_signed}. That a
gambler's optimal log-growth with side information equals the mutual information is classical
\citep{kelly1956,cover2006}, with the financial-value-of-information bound due to
\citet{barroncover1988}; the log-loss as negative log-wealth is the prequential and game-theoretic
tradition \citep{dawid1984,shafervovk2019}. The growth-rate-optimal e-variable and its identity with
a KL are from safe testing \citep{gruenwald2024} and the e-value calculus \citep{vovkwang2021};
conformal test martingales for exchangeability are \citet{vng2003}. Sequential calibration testing
by betting is closest on two fronts: \citet{arnold2023} build e-processes for probabilistic (PIT,
marginal) calibration, and \citet{shaer2023} test conditional independence by betting --- the
mechanics of (i) above --- but neither identifies the growth rate with $I(R;X)$ nor frames it as the
conformal rent. Dependence measures as conditional-coverage diagnostics are \citet{feldman2021} and
\citet{braun2025}; the information-theoretic view of prediction sets is \citet{correia2024}, which
bounds conditional entropy by set size rather than reading a gap from a calibration statistic. The
parimutuel mechanism, its log-optimal incentive, and the latent recovery it supports are
\citet{cotton_micro2022,cotton_horserace}; the continuous density conventions are
\citet{micro_density}. The contribution here is the assembly: a self-contained parimutuel derivation
of $I(R;X)$ in the conformal idiom, its identification with the rent in a deployed market, and the
pairing of a certified lower bound with an anytime-valid growth-rate estimate.

\begin{thebibliography}{99}
\bibitem[Arnold et~al.(2023)]{arnold2023} Arnold, S., Henzi, A. and Ziegel, J.~F. (2023).
Sequentially valid tests for forecast calibration. \emph{The Annals of Applied Statistics}
17(3):1909--1935. arXiv:2109.11761.
\bibitem[Barron and Cover(1988)]{barroncover1988} Barron, A.~R. and Cover, T.~M. (1988). A bound on
the financial value of information. \emph{IEEE Transactions on Information Theory} 34(5):1097--1100.
\bibitem[Braun et~al.(2025)]{braun2025} Braun, Holzm\"uller, Jordan and Bach (2025). Conditional
coverage diagnostics for conformal prediction. Preprint, arXiv:2512.11779.
\bibitem[Correia et~al.(2024)]{correia2024} Correia, A., Massoli, F.~V., Louizos, C. and Behboodi,
A. (2024). An information theoretic perspective on conformal prediction. \emph{Advances in Neural
Information Processing Systems} 37. arXiv:2405.02140.
\bibitem[Cotton()]{cotton_marginally} Cotton, P. Marginally Useful: Formalizing the Information Gap
in Conformal Prediction. Companion paper.
\bibitem[Cotton()]{cotton_signed} Cotton, P. A Feynman--Wigner-Style Diagnostic for the Efficacy of
Conformal Prediction via Signed de Finetti Representations. Companion note.
\bibitem[Cotton(2021)]{cotton_horserace} Cotton, P. (2021). Inferring relative ability from winning
probability in multi-entrant contests. doi:10.1137/19M1276261.
\bibitem[Cotton(2022)]{cotton_micro2022} Cotton, P. (2022). \emph{Microprediction: Building an Open
AI Network}. MIT Press.
\bibitem[Cover and Thomas(2006)]{cover2006} Cover, T.~M. and Thomas, J.~A. (2006). \emph{Elements of
Information Theory}, 2nd ed. Wiley. (Chapter 6, gambling and side information; Pinsker's
inequality.)
\bibitem[Dawid(1984)]{dawid1984} Dawid, A.~P. (1984). Statistical theory: the prequential approach.
\emph{Journal of the Royal Statistical Society A} 147(2):278--292.
\bibitem[Foygel Barber et~al.(2021)]{foygelbarber2021} Foygel Barber, R., Cand\`es, E.~J., Ramdas,
A. and Tibshirani, R.~J. (2021). The limits of distribution-free conditional predictive inference.
\emph{Information and Inference} 10(2):455--482. arXiv:1903.04684.
\bibitem[Gr\"unwald et~al.(2024)]{gruenwald2024} Gr\"unwald, P., de Heide, R. and Koolen, W.~M.
(2024). Safe testing. \emph{Journal of the Royal Statistical Society B} 86(5):1091--1128.
arXiv:1906.07801.
\bibitem[Kelly(1956)]{kelly1956} Kelly, J.~L. (1956). A new interpretation of information rate.
\emph{Bell System Technical Journal} 35(4):917--926.
\bibitem[Lei and Wasserman(2014)]{leiwasserman2014} Lei, J. and Wasserman, L. (2014).
Distribution-free prediction bands for non-parametric regression. \emph{Journal of the Royal
Statistical Society B} 76(1):71--96.
\bibitem[microprediction/density()]{micro_density} Cotton, P. \texttt{density}: conventions for
densities used in various games. \url{https://github.com/microprediction/density}.
\bibitem[Sejdinovic et~al.(2013)]{sejdinovic2013} Sejdinovic, D., Sriperumbudur, B., Gretton, A. and
Fukumizu, K. (2013). Equivalence of distance-based and RKHS-based statistics in hypothesis testing.
\emph{The Annals of Statistics} 41(5):2263--2291.
\bibitem[Shaer et~al.(2023)]{shaer2023} Shaer, S., Maman, G. and Aharoni, E. (2023). Model-X
sequential testing for conditional independence via testing by betting. \emph{AISTATS}.
arXiv:2210.00354.
\bibitem[Shafer and Vovk(2019)]{shafervovk2019} Shafer, G. and Vovk, V. (2019). \emph{Game-Theoretic
Foundations for Probability and Finance}. Wiley.
\bibitem[Simon-Gabriel et~al.(2023)]{simongabriel2023} Simon-Gabriel, C.-J., Barp, A., Sch\"olkopf,
B. and Mackey, L. (2023). Metrizing weak convergence with maximum mean discrepancies. \emph{Journal
of Machine Learning Research} 24. arXiv:2006.09268.
\bibitem[Sz\'ekely et~al.(2007)]{szekely2007} Sz\'ekely, G.~J., Rizzo, M.~L. and Bakirov, N.~K.
(2007). Measuring and testing dependence by correlation of distances. \emph{The Annals of
Statistics} 35(6):2769--2794.
\bibitem[Vovk(2012)]{vovk2012} Vovk, V. (2012). Conditional validity of inductive conformal
predictors. \emph{Asian Conference on Machine Learning}, PMLR 25:475--490.
\bibitem[Vovk et~al.(2003)]{vng2003} Vovk, V., Nouretdinov, I. and Gammerman, A. (2003). Testing
exchangeability on-line. \emph{International Conference on Machine Learning}.
\bibitem[Vovk and Wang(2021)]{vovkwang2021} Vovk, V. and Wang, R. (2021). E-values: calibration,
combination, and applications. \emph{The Annals of Statistics} 49(3):1736--1754.
\bibitem[Wang and Tay(2022)]{wangtay2022} Wang, C.~X. and Tay, W.~P. (2022). Practical bounds of
Kullback--Leibler divergence using maximum mean discrepancy. Preprint, arXiv:2204.02031.
\end{thebibliography}

\end{document}
